\begin{pblock}{MAREA: the archive measures its own tide}
Every published method assumes \term{one uniform water level per scene},
from an ocean model at the mouth. Inside an estuary the tide arrives late
and deformed — minutes of clock error become tens of centimetres of level
assigned to every scene. MAREA (\emph{Mudflat Altimetry via
Remotely-Estimated tides from the Archive}) measures that interior tide
from the imagery itself.

\step{1}{Each pixel is a binary threshold sensor: wet when the water tops
its elevation. Its wet/dry sequence over 465 dated scenes records
\emph{when} the water crossed it.}
\step{2}{Thousands of sensors share the same water. Band the estuary by
geodesic distance $s$ from the mouth (equal pixel counts) and ask each
band: \emph{which clock makes your wet/dry record most probable?}}
\step{3}{Score a candidate lag $\tau$ by the profiled Bernoulli
likelihood — every pixel granted its best elevation and width, so only
the clocks remain as unknowns, anchored to $\tau=0$ at the mouth:}

\vspace{-2mm}
\begin{equation*}
  \mathcal{L}(\tau)=\sum_p \max_{z_p,\sigma_p}\sum_t
  \log\!\left[P_{pt}^{\,y_{pt}}\,(1-P_{pt})^{1-y_{pt}}\right],
  \quad
  P_{pt}=\Phi\!\left(\tfrac{h_B(t-\tau)-z_p}{\sigma_p}\right)
\end{equation*}

\vspace{2mm}
\pfig[\linewidth]{m2_gate}
\figcap{The identifiability gate: a lag profile planted in a calibrated
digital twin (real sampling, real clouds) is recovered with
\emphnum{4.7\,min} RMSE, beating the literature baseline (5.4). Right: the
gain profile is machine-flat — the affine theorem, executable.}
\end{pblock}
